%!TEX root = ../main.tex

% Main Text
\section{Conclusions}

\noindent Vision-based crop-row guidance is dominated by a single recipe: a detector localises row segments as bounding boxes, a guidance line is fit through the box centroids, and the line is handed to a steering controller~\cite{diao_yolox,cth1}. Two gaps persist across this literature and across the alternative segmentation, row-anchor and curve-regression paradigms~\cite{cth1,alnet,rowdetr}. First, the systems are routinely tuned and reported in detection units (mAP), whereas the deployed quantity is a \emph{navigation} error---heading in degrees and cross-track distance in centimetres---and the two are not interchangeable. Second, regardless of paradigm, the guidance line is fit with equal or heuristic point weights, even though the far-field detections that anchor the look-ahead heading are precisely the least reliable points in the frame. Box-regression uncertainty has been modelled before, but for better \emph{detection} (KL-Loss~\cite{klloss}, Gaussian-YOLO~\cite{gaussianyolo}) rather than read out as a calibrated per-detection covariance and propagated into the line fit.

\noindent This paper has addressed both gaps with CALM-Row, a post-head module for a stock lightweight detector. CALM-Row exploits an uncertainty signal that the detector already produces for free: the Distribution Focal Loss (DFL) head~\cite{gfl} emits, per anchor, a discrete probability distribution $p_e(k)$ over $R$ bins for each of the four box edges $e\in\{l,t,r,b\}$. We read the variance $\sigma_e^2=\sum_k (k-\mu_e)^2\,p_e(k)$ of that existing distribution, propagate it by linear error propagation into a per-box centroid variance, $\sigma_{c_x}^2=(s^2/4)(\sigma_l^2+\sigma_r^2)$ (with stride $s$), and use the resulting precision $w_i=1/\sigma_{c_x,i}^2$ to weight a differentiable generalized-least-squares (GLS) fit of the guidance line $x=a\,y+b$. The heading $\theta=\arctan a$ and its uncertainty follow in closed form from the weighted parameter covariance. The detector is trained with two auxiliary losses computed entirely from this distribution and from matched ground-truth centroids, and introducing no new learnable parameters: a distance-aware calibration loss (DDC), a Gaussian negative-log-likelihood that ties the predicted centroid variance to the realised squared centroid error so that the variance becomes a \emph{calibrated} uncertainty; and a line-stability loss (LSL), whose gradients flow through the weights $w_i$ into the DFL logits so the network learns to down-weight unreliable far boxes. The ground-truth line used in LSL is fit (unweighted) through ground-truth box centroids and is therefore detector-independent, while the evaluation line is mask-derived and independent, so the central weighted-versus-unweighted claim is not evaluated circularly.

\noindent CALM-Row is deliberately \emph{not} an architectural contribution. It introduces no new operator on the feature maps and no new learnable parameters in the training losses; all computation is post-head, CPU-side arithmetic over post-NMS boxes. The exported INT8/RKNN (Rockchip) or TensorRT (Jetson) graph is therefore a vanilla YOLO with zero added NPU operations, and because the module reads only the standard DFL distribution it is detector-agnostic, applying without modification to YOLOv8n, YOLOv10n and any DFL-headed successor. A synthetic mechanism check, clearly labelled as illustrative, is reported in Section~\ref{sec:res-synth}.

\noindent We have specified, rather than yet executed, an evaluation designed around the deployed task. The primary metric is controller-independent heading error in degrees, with line-fit RMSE (px and cm via a stated ground-plane homography) and a static geometric cross-track proxy in centimetres as secondary measures (a closed-loop pure-pursuit simulation is left to future work); calibration is reported via expected calibration error and reliability, before and after INT8 quantization. The core comparison is a paired, per-frame A/B between equal-weight least squares (the de-facto recipe), a readout-only arm that applies the raw DFL variance with no extra training, and the full CALM-Row, isolating the gain from weighting alone versus from calibration-aware training. Generalization is assessed in-domain on the CRDLD~\cite{cth1} test split and as cross-dataset transfer to CRBD~\cite{crbd} (different camera, crop and resolution) without retraining, with reliability stratified continuously against an image-row range proxy rather than categorical distance bins. All confirmatory comparisons use multiple seeds, paired non-parametric tests with multiple-comparison correction, and effect sizes with bootstrap confidence intervals; published crop-row systems~\cite{diao_yolox,rowdetr,alnet} are positioned on the same test split, ground-truth line and edge device. We deliberately make no quantitative claims here: all navigation, calibration and efficiency numbers remain to be measured under this protocol [TODO: report measured values once experiments are run].

\noindent Several directions extend this work. The most important is real closed-loop validation: replacing the simulator with field trials on the named edge platforms, where measured perception latency, controller dynamics and terrain interact with the guidance estimate. A second direction is temporal consistency---fusing per-frame line estimates and their calibrated heading variance across time, for example with an uncertainty-aware filter, to suppress frame-to-frame jitter at the look-ahead row. A third is learned geometric range supervision: where stereo, LiDAR or known inter-row spacing is available, the currently free, monocular uncertainty could be tied to true metric range and ground-plane geometry, turning the calibrated covariance from an image-space proxy into a metrically grounded reliability estimate without compromising the zero-NPU-cost, detector-agnostic deployment that motivates CALM-Row.